\documentclass[11pt,a4paper,addpoints]{exam}
\usepackage[utf8]{inputenc}
\usepackage[T1]{fontenc}
\usepackage[english]{babel}
\usepackage[margin=1.8cm]{geometry}

\pointpoints{pts}{pts}
\boxedpoints

\lhead{\makebox[12cm][l]{\large\textbf{Name and surname:}\enspace\hrulefill}}
\rhead{W04: collections-generics}

\footer{}{}{}

\renewcommand*\half{.5}

\begin{document}
\emergencystretch=2em

\cellwidth{2.75cm}
\noindent \gradetable[h]

\vspace{0.4cm}

\noindent Gustavo Fring from Los Pollos Hermanos needs your help in his restaurant. You are provided with a pre-configured solution \texttt{LosPollosHermanos.slnx} containing the application project \texttt{LosPollosHermanos} and the unit test project \texttt{LosPollosHermanos.Tests}. The files \texttt{Program.cs} and \texttt{Records.cs} are already provided and must not be modified. The active stage is defined in the \texttt{LosPollosHermanos.props} file visible in the solution (conditional compilation symbols \texttt{STAGE1}, \texttt{STAGE2}, \texttt{STAGE3}), which is imported by both projects. Uncomment the corresponding symbol before starting each stage. \textbf{Stage 1} must be completed before \textbf{Stage 2}; \textbf{Stage 3} is independent. Are you up to Los Pollos standards?

\begin{questions}

\question[1\half] In file \texttt{Line.cs}, implement the following methods:
\begin{itemize}
    \item \texttt{Add} -- adds an item at the \textbf{back} of the line, returns \texttt{false} if full,
    \item \texttt{Take} -- removes and returns the element from the \textbf{front} of the line, returns \texttt{null} if empty,
    \item \texttt{Check} -- returns the element from the \textbf{front} of the line without removing it, returns \texttt{null} if empty,
    \item \texttt{GetEnumerator} -- returns an iterator yielding elements from the line \textbf{front to back} (FIFO).
\end{itemize}

\question[2\half] In file \texttt{Kitchen.cs}, implement the following methods:
\begin{itemize}
    \item \texttt{StockEquipment} -- adds new equipment to \texttt{\_equipment} if not already present and returns how many items were newly added,
    \item \texttt{StockSpices} -- adds the given amounts of spices to \texttt{\_spices} stock, ensuring each spice does not exceed \texttt{SpiceLimit}, and returns how many grams were actually added \textbf{in total},
    \item \texttt{CanCook} -- returns \texttt{true} \textbf{if} the kitchen has all required equipment (\texttt{recipe.Equipment}) and enough of every spice listed in \texttt{recipe.SpicesPerBatch},
    \item \texttt{CookBatches} -- for \texttt{attempt = 0, 1, \dots, count - 1} (\textbf{at most} \texttt{count} attempts), in order:
    \begin{itemize}
        \item checks via \texttt{CanCook} if the recipe \textbf{can be cooked}; if not, \textbf{stops} the cook,
        \item consumes the spices required by the recipe (\texttt{recipe.SpicesPerBatch}),
        \item builds a \texttt{ChickenBatch} (\texttt{id + attempt}, \texttt{pieces}, \texttt{recipe.Temperature}, and \texttt{HasMeth} set to \texttt{true} when \texttt{attempt \% methEvery == 0}),
        \item checks if the batch meets Los Pollos standards via \texttt{Accept(batch, attempt)}; if not, \textbf{stops} the cook,
        \item adds the batch to the \texttt{\_ready} line; if full, \textbf{stops} the cook,
        \item only then \textbf{yields} the batch.
    \end{itemize}
\end{itemize}

\question[1]
In file \texttt{BatchComparer.cs}, complete \texttt{BatchComparer} (\texttt{Comparer<ChickenBatch>}) by implementing the \texttt{Compare} method using the following ordering key:
    \begin{enumerate}
        \item safe batches before unsafe (a batch is safe if \texttt{Temperature >= SafeTemperature}),
        \item batches containing meth (\texttt{HasMeth == true}) before clean ones,
        \item smaller batches before bigger -- ascending by \texttt{Pieces}.
    \end{enumerate}
\textbf{Remember to handle} \texttt{null}\textbf{s}: \texttt{null} sorts last, and comparing two \texttt{null} references returns \texttt{0}.
\end{questions}

\end{document}
